\documentclass[11pt]{article}

\usepackage[margin=0.9in]{geometry}
\usepackage{amsmath,amssymb}
\usepackage{booktabs}
\usepackage{enumitem}
\usepackage{listings}
\usepackage{xcolor}
\usepackage{hyperref}
\usepackage{tikz}

\definecolor{backcolour}{rgb}{0.95,0.95,0.92}
\definecolor{codeblue}{rgb}{0.05,0.15,0.55}
\definecolor{codegreen}{rgb}{0.0,0.45,0.0}

\lstset{
    backgroundcolor=\color{backcolour},
    basicstyle=\ttfamily\small,
    keywordstyle=\color{codeblue}\bfseries,
    commentstyle=\color{codegreen},
    stringstyle=\color{purple},
    breaklines=true,
    frame=single,
    showstringspaces=false,
    tabsize=4
}

\title{CPSC 250 \\ Plotting the Credits Profile}
\author{Edward Brash}
\date{}

\begin{document}

\maketitle

\section{Goal}

After simulating many games, we want to visualize the results.

In particular, we want to plot the number of credits as a function of round number.

This is a data manipulation task:

\begin{quote}
Store the data, then plot the data.
\end{quote}

\section{What Data Do We Need?}

To plot credits versus round number, we need two lists:

\begin{itemize}
    \item x-values: round numbers,
    \item y-values: credits.
\end{itemize}

For example:

\begin{lstlisting}[language=Python]
rounds = [0, 1, 2, 3, 4]
credits = [10, 11, 12, 11, 12]
\end{lstlisting}

These lists must correspond element-by-element.

\section{Storing Simulation Results}

For one simulation, we might store:

\begin{lstlisting}[language=Python]
result_round = []
credit_round = []
\end{lstlisting}

At each round:

\begin{lstlisting}[language=Python]
result_round.append(rounds)
credit_round.append(credits)
\end{lstlisting}

For many simulations, we need lists of lists:

\begin{lstlisting}[language=Python]
round_result = []
credit_result = []
result_number = []
\end{lstlisting}

Each simulation contributes one list of round numbers and one list of credit values.

\section{Matplotlib}

Python's most common plotting library is Matplotlib.

The usual import is:

\begin{lstlisting}[language=Python]
import matplotlib.pyplot as plt
\end{lstlisting}

The \texttt{pyplot} part is one submodule of Matplotlib.

By convention, we import it as:

\begin{lstlisting}[language=Python]
plt
\end{lstlisting}

\section{Basic Plot}

A basic line plot is:

\begin{lstlisting}[language=Python]
plt.plot(result_round, credit_round)
plt.show()
\end{lstlisting}

Here:

\begin{itemize}
    \item \verb|result_round| provides the x-values,
    \item \verb|credit_round| provides the y-values.
\end{itemize}

\section{Rules of Plotting}

A good plot should include:

\begin{enumerate}
    \item a descriptive title,
    \item a labeled x-axis,
    \item a labeled y-axis.
\end{enumerate}

For this problem, a good title is:

\begin{lstlisting}[language=Python]
plt.title("Credits vs. Round Number")
\end{lstlisting}

The axis labels are:

\begin{lstlisting}[language=Python]
plt.xlabel("Round Number")
plt.ylabel("Credits")
\end{lstlisting}

\section{Complete Single-Game Plot}

\begin{lstlisting}[language=Python]
plt.plot(result_round, credit_round)

plt.title("Credits vs. Round Number")
plt.xlabel("Round Number")
plt.ylabel("Credits")

plt.show()
\end{lstlisting}

\section{Plotting Multiple Simulations}

If we run many simulations, we can plot several credit histories at once.

Suppose:

\begin{lstlisting}[language=Python]
round_result[i]
credit_result[i]
\end{lstlisting}

store the data for simulation \(i\).

Then:

\begin{lstlisting}[language=Python]
for i in range(num_sims):
    plt.plot(round_result[i], credit_result[i])
\end{lstlisting}

plots all simulations on the same axes.

\section{Interpreting the Plot}

Each line represents one simulated game.

A game begins with some initial number of credits.

Then credits move up and down until the game ends at zero credits.

This creates a random-walk-like path.

\section{Histogram of Game Lengths}

Another useful plot is a histogram of the number of rounds required before the game ends.

If:

\begin{lstlisting}[language=Python]
result_number = [18422, 503, 971, ...]
\end{lstlisting}

stores the number of rounds for each simulation, then:

\begin{lstlisting}[language=Python]
plt.hist(result_number)
plt.show()
\end{lstlisting}

creates a histogram.

\section{Number of Bins}

A histogram groups data into bins.

The number of bins affects the visual interpretation.

For example:

\begin{lstlisting}[language=Python]
plt.hist(result_number, bins=20)
\end{lstlisting}

uses 20 bins.

Too few bins may hide structure. Too many bins may make the plot noisy.

\section{Histogram Labels}

As always, we label the plot.

\begin{lstlisting}[language=Python]
plt.title("Distribution of Game Lengths")
plt.xlabel("Number of Rounds")
plt.ylabel("Frequency")
\end{lstlisting}

\section{Complete Histogram Example}

\begin{lstlisting}[language=Python]
plt.hist(result_number, bins=20)

plt.title("Distribution of Game Lengths")
plt.xlabel("Number of Rounds")
plt.ylabel("Frequency")

plt.show()
\end{lstlisting}

\section{Philosophy of the Simulation}

The notes emphasize that we are not just playing a game.

We are modeling a random walk process.

This connects programming to probability, statistics, and simulation.

The general scientific workflow is:

\begin{enumerate}
    \item define a model,
    \item run the model many times,
    \item store the results,
    \item visualize the results,
    \item interpret the behavior.
\end{enumerate}

\section{Key Takeaways}

\begin{itemize}
    \item Plotting requires stored x-values and y-values.
    \item For the credits simulation, x is round number and y is credits.
    \item Matplotlib is commonly imported as \verb|plt|.
    \item Good plots need titles and axis labels.
    \item Multiple simulations can be plotted on the same axes.
    \item Histograms show distributions.
    \item The number of bins controls histogram resolution.
\end{itemize}

\end{document}
